\documentclass{fiche}
\metadonnees{1}{Portraits}{Bloc 1 — Le récit : présent et passé}

\begin{document}
\entete

\objectifs{%
\textbf{Langue} : présent simple et présent en \emph{be + -ing} ; verbes d'état ; lexique du portrait.\par
\textbf{Texte} : Louisa May Alcott, \emph{Little Women} (1868, texte de l'édition revue de 1880), chapitre \textsc{i}.\par
\textbf{Durée} : 1\,h\,30 — exercices 35 min, texte et questions 30 min, version 25 min.}

%% ===================================================================
\exercice[12 min]{Présent simple ou présent en \emph{be + -ing}}

\rappel{\textbf{Rappel.} Le présent simple dit ce qui est vrai en général, habituel ou
permanent. Le présent en \emph{be + -ing} dit ce qui est en cours au moment où l'on parle,
ou ce qui est provisoire. Repères du présent simple : \emph{usually}, \emph{every day},
\emph{never} ; repères du \emph{be + -ing} : \emph{look}, \emph{listen}, \emph{now},
\emph{at the moment}, \emph{this winter}. Les verbes d'état n'ont pas de forme en
\emph{be + -ing}.}

\consigne{Mettez le verbe entre parenthèses à la forme qui convient.}

\begin{anglais}
\begin{enumerate}[label=\arabic*.,itemsep=2.2mm,leftmargin=*]
  \item Look! Jo \emph{(lie)} \trou{is lying} on the rug again.
  \item Meg \emph{(teach)} \trou{teaches} the King children every morning, from nine to one.
  \item ``Why \emph{(you / whistle)} \trou{are you whistling}?'' asked Amy. ``Because it annoys you.''
  \item The March girls usually \emph{(spend)} \trou{spend} their evenings by the fire.
  \item Listen --- Beth \emph{(play)} \trou{is playing} the old piano in the parlour.
  \item Their father \emph{(not fight)} \trou{is not fighting}; he \emph{(serve)} \trou{is serving} as a chaplain this winter.
  \item Amy \emph{(always / use)} \trou{is always using} long words she does not understand.
  \item Water \emph{(freeze)} \trou{freezes} at zero degrees.
  \item This December the snow \emph{(fall)} \trou{is falling} earlier than usual.
  \item What \emph{(Jo / think)} \trou{does Jo think} of growing up?
  \item Beth \emph{(think)} \trou{is thinking} of a present for her mother at this very moment.
  \item Meg \emph{(not like)} \trou{does not like} her old dress.
\end{enumerate}
\end{anglais}

\note{Présent simple : 2, 4, 8, 12. Présent en \emph{be + -ing} : 1, 3, 5, 6, 9, 11.}

\note[Deux cas particuliers]{Phrase 7 : \emph{always + be + -ing} n'exprime pas une action en
cours mais l'agacement du locuteur (\og elle n'arrête pas d'employer\ldots \fg). Phrases 10
et 11 : \emph{think} au sens de \og avoir une opinion \fg{} est un verbe d'état (présent
simple) ; au sens de \og réfléchir, songer à \fg{} il devient un verbe d'action
(\emph{be + -ing}).}

%% ===================================================================
\exercice[10 min]{Verbes d'état}

\rappel{\textbf{Rappel.} Un verbe d'état décrit au lieu de raconter : \emph{know},
\emph{believe}, \emph{understand}, \emph{want}, \emph{like}, \emph{hate}, \emph{need},
\emph{see}, \emph{hear}, \emph{taste}, \emph{seem}, \emph{have} au sens de \og posséder \fg.
Il n'a pas de forme en \emph{be + -ing}. Mais plusieurs de ces verbes changent de sens et
deviennent des verbes d'action : \emph{to have lunch} (prendre), \emph{to see somebody}
(fréquenter), \emph{to taste the soup} (goûter).}

\consigne{Six de ces phrases contiennent une erreur ; deux sont correctes. Corrigez les
formes fautives et cochez les phrases correctes.}

\begin{anglais}
\begin{enumerate}[label=\arabic*.,itemsep=2.2mm,leftmargin=*]
  \item I am not believing a word of it. \hfill \trou{I do not believe}
  \item Amy is having a new box of drawing pencils. \hfill \trou{Amy has}
  \item They are having tea in the parlour. \hfill \trou{correct}
  \item Meg is wanting a pair of gloves for her mother. \hfill \trou{Meg wants}
  \item I am knowing the answer. \hfill \trou{I know}
  \item Beth is being very quiet this evening. \hfill \trou{correct}
  \item Jo is seeing the point now. \hfill \trou{Jo sees}
  \item This soup is tasting of nothing. \hfill \trou{tastes}
\end{enumerate}
\end{anglais}

\note[Les deux phrases correctes]{3 : \emph{have} n'y signifie pas \og posséder \fg{} mais
\og prendre \fg{} — c'est alors un verbe d'action. 6 : \emph{be + -ing} avec un adjectif
décrit un comportement passager (\og Beth est bien silencieuse ce soir \fg) et non une
qualité permanente (\emph{Beth is quiet} = elle est de nature silencieuse).}

%% ===================================================================
\exercice[13 min]{Le portrait : vocabulaire}
\consigne{a) Associez chaque mot à sa traduction. Ces mots reviennent tous dans le texte.}

\begin{multicols}{2}
\begin{enumerate}[label=\arabic*.,itemsep=1.2mm,leftmargin=*]
  \item plump \hfill \trou{f}
  \item fair \hfill \trou{i}
  \item slender \hfill \trou{a}
  \item shy \hfill \trou{k}
  \item timid \hfill \trou{c}
  \item vain \hfill \trou{h}
  \item rosy \hfill \trou{d}
  \item thoughtful \hfill \trou{l}
  \item fierce \hfill \trou{b}
  \item prim \hfill \trou{g}
  \item affected \hfill \trou{j}
  \item decided \hfill \trou{e}
\end{enumerate}
\columnbreak
\begin{enumerate}[label=\alph*.,itemsep=1.2mm,leftmargin=*]
  \item svelte, élancé
  \item farouche, féroce
  \item craintif
  \item au teint frais
  \item résolu, décidé
  \item potelé, dodu
  \item guindé, collet monté
  \item vaniteux
  \item clair, blond
  \item maniéré, affecté
  \item timide, réservé
  \item pensif, songeur
\end{enumerate}
\end{multicols}

\note[Faux amis]{\emph{fair} ne signifie ici ni \og juste \fg{} ni \og foire \fg{} mais
\og clair \fg{} (teint, cheveux). \emph{affected} n'est pas \og affecté, touché \fg{} mais
\og maniéré \fg. \emph{decided} appliqué à une bouche : \og une bouche résolue \fg.}

\consigne{b) Décrivez en cinq phrases une personne que vous connaissez. Employez au moins
quatre mots de la liste, un présent en \emph{be + -ing} et un verbe d'état.}
\lignes{6}
\corr{Réponse libre. Attendus : accord du présent simple à la 3\textsuperscript{e} personne
(\emph{she looks}, \emph{he seems}) ; un seul \emph{be + -ing} justifié par un repère
temporel ; aucun \emph{be + -ing} sur un verbe d'état.}

%% ===================================================================
\newpage
\section{Texte}

\noindent\small\textit{Il est question des quatre sœurs March, en Nouvelle-Angleterre,
pendant la guerre de Sécession. Leur père est parti au front.}\normalsize

\begin{texte}
``Christmas won't be Christmas without any presents,'' grumbled Jo, lying on the rug.

``It's so dreadful to be poor!'' sighed Meg, looking down at her old dress.

``I don't think it's fair for some girls to have plenty of pretty things, and other girls
nothing at all,'' added little Amy, with an \gl[an injured sniff]{injured sniff}{un
reniflement offensé}.

``We've got father and mother and each other,'' said Beth contentedly, from her corner.

The four young faces on which the firelight shone brightened at the cheerful words, but
darkened again as Jo said sadly, ``We haven't got father, and shall not have him for a long
time.'' She didn't say ``perhaps never,'' but each silently added it, thinking of father far
away, where the fighting was.

Nobody spoke for a minute; then Meg said in an altered tone, ``You know the reason mother
proposed not having any presents this Christmas was because it is going to be a hard winter
for every one; and she thinks we ought not to spend money for pleasure, when our men are
suffering so in the army. We can't do much, but we can make our little sacrifices, and ought
to do it gladly. But I am afraid I don't;'' and Meg shook her head, as she thought
regretfully of all the pretty things she wanted.

``But I don't think the little we should spend would do any good. We've each got a dollar,
and the army wouldn't be much helped by our giving that. I agree not to expect anything from
mother or you, but I do want to buy \gl[Undine and Sintram]{\emph{Undine and
Sintram}}{roman allemand de La Motte-Fouqué, 1811} for myself; I've wanted it \emph{so}
long,'' said Jo, who was a bookworm.

``I planned to spend mine in new music,'' said Beth, with a little sigh, which no one heard
but the \gl[the hearth-brush]{hearth-brush}{la brosse de cheminée} and
\gl[a kettle-holder]{kettle-holder}{une manique}.

``I shall get a nice box of Faber's drawing-pencils; I really need them,'' said Amy
decidedly.

``Mother didn't say anything about our money, and she won't wish us to give up everything.
Let's each buy what we want, and have a little fun; I'm sure we work hard enough to earn
it,'' cried Jo, examining the heels of her shoes in a gentlemanly manner.

``I know \emph{I} do, --- teaching those \gl[tiresome]{tiresome}{pénible, assommant} children
nearly all day, when I'm longing to enjoy myself at home,'' began Meg, in the complaining
tone again.

``You don't have half such a hard time as I do,'' said Jo. ``How would you like to be shut up
for hours with a nervous, fussy old lady, who \gl[to keep somebody trotting]{keeps you
trotting}{faire trotter quelqu'un sans arrêt}, is never satisfied, and worries you till
you're ready to fly out of the window or cry?''

``It's naughty to fret; but I do think washing dishes and keeping things tidy is the worst
work in the world. It makes me cross; and my hands get so \gl[stiff]{stiff}{raide,
engourdi}, I can't practise well at all;'' and Beth looked at her rough hands with a sigh
that any one could hear that time.

``I don't believe any of you suffer as I do,'' cried Amy; ``for you don't have to go to
school with impertinent girls, who \gl[to plague somebody]{plague}{harceler quelqu'un} you if
you don't know your lessons, and laugh at your dresses, and label your father if he isn't
rich, and insult you when your nose isn't nice.''

``If you mean \gl[libel]{\emph{libel}}{la diffamation ; Amy a dit \emph{label},
\og étiqueter \fg}, I'd say so, and not talk about \emph{labels}, as if papa was a
\gl[a pickle-bottle]{pickle-bottle}{un bocal de cornichons},'' advised Jo, laughing.

``I know what I mean, and you needn't be \emph{statirical} about it. It's proper to use good
words, and improve your \emph{vocabilary},'' returned Amy, with dignity.

``Don't \gl[to peck at one another]{peck at one another}{se chamailler, litt. se donner des
coups de bec}, children. Don't you wish we had the money papa lost when we were little, Jo?
Dear me! how happy and good we'd be, if we had no worries!'' said Meg, who could remember
better times.

``You said the other day you thought we were \gl[a deal happier]{a deal happier}{bien plus
heureux (tour familier)} than the King children, for they were fighting and fretting all the
time, in spite of their money.''

``So I did, Beth. Well, I think we are; for, though we do have to work, we make fun for
ourselves, and are a pretty jolly set, as Jo would say.''

``Jo does use such slang words!'' observed Amy, with a reproving look at the long figure
stretched on the rug. Jo immediately sat up, put her hands in her pockets, and began to
whistle.

``Don't, Jo; it's so boyish!''

``That's why I do it.''

``I detest rude, \gl[unlady-like]{unlady-like}{peu distingué} girls!''

``I hate affected, \gl[a niminy-piminy chit]{niminy-piminy chits}{une pimbêche minaudière}!''

``\,`Birds in their little nests agree,'\,'' sang Beth, the peace-maker, with such a funny
face that both sharp voices softened to a laugh, and the ``pecking'' ended for that time.
\end{texte}
\source{Louisa May Alcott, \emph{Little Women}, ch.~\textsc{i} (texte de l'édition revue,
Boston, Roberts Brothers, 1880).}

\partiel[15 min]{Questions}
\consigne{\en{Answer in English, in full sentences.}}

\begin{enumerate}[label=\arabic*.,itemsep=2mm,leftmargin=*]
  \item \en{Why are the girls going to have no presents this Christmas?}
    \lignes{2}
    \corr{Their mother suggested it because the winter will be hard for everyone and she
    thinks they should not spend money on pleasure while the soldiers are suffering.}
  \item \en{What do we learn about their father in this passage?}
    \lignes{2}
    \corr{He is far away, where the fighting is, and will not be home for a long time. The
    girls fear he may never come back, but none of them says so aloud.}
  \item \en{Each sister complains about something. Say what, in one sentence for each.}
    \lignes{4}
    \corr{Meg is tired of teaching the King children all day. Jo hates being shut up with a
    fussy old lady who is never satisfied. Beth finds washing dishes the worst work in the
    world, and it makes her hands stiff. Amy is laughed at by impertinent girls at school.}
  \item \en{Amy says \emph{label} instead of \emph{libel}. What does the mistake tell us
    about her, and how do the others react?}
    \lignes{3}
    \corr{She wants to sound grown-up and elegant, and uses long words she does not master.
    Jo laughs at her and mocks her (\emph{as if papa was a pickle-bottle}), while Amy defends
    herself with dignity and makes two more mistakes (\emph{statirical}, \emph{vocabilary}).}
  \item \en{Choose two adjectives that fit Jo and justify each one with a short quotation.}
    \lignes{3}
    \corr{Par exemple \emph{boyish} (\og put her hands in her pockets, and began to
    whistle \fg), \emph{outspoken} ou \emph{rebellious} (\og That's why I do it \fg),
    \emph{bookish} (\og said Jo, who was a bookworm \fg).}
  \item \en{Find three \emph{-ing} forms in the text. For each, say whether it is a
    present continuous or something else.}
    \lignes{3}
    \corr{\emph{lying on the rug}, \emph{looking down at her old dress},
    \emph{examining the heels of her shoes} : participes présents apposés, qui décrivent
    une action simultanée — ce ne sont pas des présents continus (pas d'auxiliaire
    \emph{be}). Vrais présents continus : \emph{when I'm longing to enjoy myself},
    \emph{our men are suffering}. \emph{they were fighting and fretting} est un prétérit
    en \emph{be + -ing}.}
\end{enumerate}

\note{Une forme en \emph{-ing} n'est un présent continu que si elle est précédée de
\emph{am / is / are}. Sans auxiliaire, c'est un participe présent (l'équivalent du gérondif
français en \og en\ldots ant \fg) ou un nom verbal (\emph{washing dishes and keeping things
tidy is the worst work}, sujet du verbe \emph{is}).}

%% ===================================================================
\newpage
\section{Version}
\consigne{Traduisez le passage en français. Il vient quelques lignes plus loin dans le même
chapitre. La phrase d'introduction vous est donnée : \og Comme les jeunes lecteurs aiment
savoir \emph{de quoi les gens ont l'air}, nous profiterons de ce moment pour leur faire une
petite esquisse des quatre sœurs, assises à tricoter dans le crépuscule, tandis que la neige
de décembre tombait doucement dehors et que le feu crépitait gaiement dedans. \fg}

\begin{version}
Margaret, the eldest of the four, was sixteen, and very pretty, being plump and fair, with
large eyes, plenty of soft, brown hair, a sweet mouth, and white hands, of which she was
rather vain. Fifteen-year-old Jo was very tall, thin, and brown, and reminded one of a
\gl[a colt]{colt}{un poulain}; for she never seemed to know what to do with her long limbs,
which were very much in her way. She had a decided mouth, a comical nose, and sharp, gray
eyes, which appeared to see everything, and were by turns fierce, funny, or thoughtful. Her
long, thick hair was her one beauty; but it was usually bundled into a net, to be out of her
way. Round shoulders had Jo, big hands and feet, a \gl[a fly-away look]{fly-away look}{un air
flottant, négligé} to her clothes, and the uncomfortable appearance of a girl who was rapidly
shooting up into a woman, and didn't like it. Elizabeth --- or Beth, as every one called her
--- was a rosy, smooth-haired, bright-eyed girl of thirteen, with a shy manner, a timid
voice, and a peaceful expression, which was seldom disturbed. Her father called her ``Little
Tranquillity,'' and the name suited her excellently; for she seemed to live in a happy world
of her own, only venturing out to meet the few whom she trusted and loved. Amy, though the
youngest, was a most important person, --- in her own opinion at least. A regular
\gl[a snow-maiden]{snow-maiden}{une fée des neiges}, with blue eyes, and yellow hair, curling
on her shoulders, pale and slender, and always carrying herself like a young lady
\gl[mindful of]{mindful}{soucieux de} of her manners. What the characters of the four sisters
were we will leave to be found out.
\end{version}
\source{\emph{Ibid.}}

\lignes{22}

\corr{\textbf{Proposition de traduction.} Margaret, l'aînée des quatre, avait seize ans ;
elle était très jolie, potelée et le teint clair, avec de grands yeux, une abondante
chevelure brune et soyeuse, une bouche charmante et des mains blanches dont elle tirait
quelque vanité. Jo, quinze ans, était très grande, mince et hâlée, et faisait songer à un
poulain, car elle ne semblait jamais savoir que faire de ses longs membres, qui la gênaient
considérablement. Elle avait la bouche résolue, le nez comique et des yeux gris perçants qui
paraissaient tout voir et qui étaient tour à tour farouches, malicieux ou pensifs. Ses
cheveux longs et épais étaient sa seule beauté, mais elle les fourrait d'ordinaire dans une
résille pour ne pas être gênée. Jo avait les épaules rondes, de grandes mains et de grands
pieds, quelque chose de flottant dans sa mise, et l'air emprunté d'une fille qui pousse à
toute vitesse vers l'âge de femme et qui n'aime pas cela. Elizabeth --- ou Beth, comme tout
le monde l'appelait --- était une enfant de treize ans au teint frais, aux cheveux lisses et
aux yeux vifs, avec des manières timides, une voix craintive et une expression paisible que
rien ne venait presque jamais troubler. Son père l'appelait \og Petite Tranquillité \fg, et
le nom lui allait à merveille : elle semblait vivre dans un monde heureux bien à elle, dont
elle ne sortait que pour aller vers les rares êtres en qui elle avait confiance et qu'elle
aimait. Amy, quoique la plus jeune, était un personnage de la première importance --- à ses
propres yeux du moins. Une véritable fée des neiges, avec ses yeux bleus et ses cheveux
blonds bouclant sur ses épaules, pâle et svelte, et se tenant toujours comme une demoiselle
soucieuse de ses manières. Quant au caractère des quatre sœurs, nous laisserons au lecteur le
soin de le découvrir.}

\note[Choix de traduction 1]{\emph{plump and fair} : \emph{fair} porte ici sur le teint, non
sur les cheveux (ils sont \emph{brown} deux mots plus loin). De même \emph{thin and brown}
pour Jo : \emph{brown} = hâlée par le grand air, et non \og brune \fg.}

\note[Choix de traduction 2]{\emph{reminded one of a colt} : \emph{one} est le pronom
impersonnel, exactement notre \og on \fg. \og Elle rappelait à quelqu'un un poulain \fg{}
serait un contresens.}

\note[Choix de traduction 3]{\emph{a rosy, smooth-haired, bright-eyed girl} : l'anglais
entasse les adjectifs composés devant le nom ; le français ne le peut pas et doit les
déplier après le nom (\og une enfant au teint frais, aux cheveux lisses, aux yeux vifs \fg).
C'est l'étoffement, le réflexe le plus rentable de la version.}

\note[Choix de traduction 4]{\emph{a most important person} : \emph{most} sans article
défini n'est pas un superlatif mais un intensif, \og des plus importants \fg, \og de la
première importance \fg. Comparer : \emph{the most important person} = la plus importante.}

\note[Choix de traduction 5]{\emph{Round shoulders had Jo} : inversion expressive, sans
équivalent direct en français. On rétablit l'ordre normal, quitte à compenser ailleurs par
un rythme marqué (\og Jo avait les épaules rondes, de grandes mains et de grands pieds\ldots \fg).}

\note[Choix de traduction 6]{\emph{which were very much in her way} puis \emph{to be out of
her way} : l'anglais joue sur la reprise d'une même expression. Le français peut la garder
(\og qui la gênaient \fg{} / \og pour ne pas être gênée \fg) : c'est un effet, pas une
maladresse.}

\note[Temps du récit]{Toute la version est au prétérit, temps du récit. Il fait l'objet de la
fiche 2 ; ici on se contente de le rendre par l'imparfait (description) plutôt que par le
passé simple (action ponctuelle).}

%% ===================================================================
\partiel{Lexique à mémoriser}
\begin{itemize}[label=--,itemsep=0.8mm,leftmargin=*]\small
  \item \textbf{plump} — potelé
  \item \textbf{fair} — clair (teint, cheveux) ; \textit{même mot que dans} fair play \textit{: l'idée de ce qui est net, régulier}
  \item \textbf{slender} — svelte, élancé
  \item \textbf{shy} — timide, réservé ; \textit{d'où} to shy away from \textit{« se dérober devant »}
  \item \textbf{timid} — craintif ; \textit{même racine latine que} timide, \textit{mais plus près de} timoré
  \item \textbf{vain} — vaniteux ; \textit{cf.} vanity, in vain
  \item \textbf{rosy} — au teint frais ; \textit{suffixe} -y \textit{= « qui a la qualité de » :} rosy, rocky, dusty
  \item \textbf{thoughtful} — pensif ; \textit{suffixe} -ful \textit{= « plein de » :} thought + ful
  \item \textbf{fierce} — farouche ; \textit{du latin} ferus \textit{« sauvage », comme} féroce \textit{et} fier
  \item \textbf{prim} — guindé ; \textit{surtout dans} prim and proper
  \item \textbf{affected} — maniéré ; \textit{participe de} to affect \textit{au sens de « feindre »}
  \item \textbf{dreadful} — affreux ; \textit{de} to dread \textit{« redouter » +} -ful
  \item \textbf{to fret} — se tracasser ; \textit{sens premier « ronger », d'où} fretful
  \item \textbf{to grumble} — ronchonner ; \textit{formation imitative, comme} grommeler
  \item \textbf{to sigh} — soupirer
  \item \textbf{to long for} — aspirer à ; \textit{de} long \textit{: trouver le temps long}
  \item \textbf{fussy} — tatillon ; \textit{de} fuss \textit{« histoires, chichis »}
  \item \textbf{a bookworm} — un rat de bibliothèque ; \textit{littéralement « ver de livre »}
  \item \textbf{slang} — l'argot
  \item \textbf{to be cross} — être de mauvaise humeur ; \textit{de} cross \textit{au sens de « de travers »}
\end{itemize}

\end{document}
